\chapter{Marco tecnológico}

Este capítulo justifica las tecnologías con las que se construyó el sistema. No pretende describirlas ---su documentación oficial lo hace mejor--- sino explicar por qué se eligió cada una frente a las alternativas disponibles, porque una decisión tecnológica sin alternativa considerada es indistinguible de una costumbre. La selección persigue tres propiedades: coherencia entre componentes, soporte sólido para pruebas y facilidad de reproducción. No se buscó maximizar el número de herramientas: cada incorporación debía resolver un problema que las ya presentes no cubrían, y en varios puntos la decisión defendible resultó ser no añadir nada.

\section{Backend y persistencia}

El backend usa ASP.NET Core sobre .NET 10 \cite{aspnetcore}. La plataforma reúne servidor \gls{http}, inyección de dependencias, configuración, validación, autenticación, autorización y registro. Frente a ensamblar esas piezas sobre un framework mínimo de JavaScript, ofrece convenciones tipadas y un pipeline integrado adecuado para una \gls{api} con reglas, roles y pruebas. La elección tampoco obliga a una arquitectura empresarial compleja: controladores, servicios y EF Core son suficientes para el tamaño del proyecto.

La plataforma admite también F\#, pero C\# es su lenguaje por defecto y el que documentan sus bibliotecas, de modo que aquí la decisión se reduce a no apartarse del camino previsto sin una razón de dominio que lo justifique. C\# permite expresar contratos mediante tipos, registros y nulabilidad. Las peticiones y respuestas son \glspl{dto} separados de las entidades para no exponer directamente el modelo persistente. La inyección por constructor hace explícitas dependencias como el contexto, el \emph{hasher} o la identidad actual, y facilita sustituirlas en pruebas.

Entity Framework Core actúa como \gls{orm} y administra el esquema mediante \glspl{migracion} versionadas \cite{efcore}. Se prefirió frente a construir \gls{sql} manual para cada acceso porque el proyecto necesita evolución incremental, relaciones y proyecciones repetidas. Ello no elimina el diseño relacional: índices filtrados, restricciones \texttt{CHECK}, tipos temporales y políticas de borrado se configuran explícitamente. Las consultas de lectura emplean \texttt{AsNoTracking} y proyecciones para no cargar grafos innecesarios.

PostgreSQL 16 proporciona persistencia relacional, transacciones, claves foráneas e índices parciales \cite{postgresql}. Un almacén documental habría simplificado poco: el dominio contiene relaciones estables y reglas de integridad que encajan de forma natural en un esquema relacional. Los identificadores técnicos son \glspl{uuid}; los identificadores de negocio, como referencia o matrícula, conservan sus propias restricciones. Las fechas que representan instantes se almacenan en \gls{utc} mediante \texttt{timestamp with time zone}.

La elección dentro de lo relacional tampoco resulta indiferente, porque dos reglas del dominio acabaron sostenidas por el motor y no por el servicio: la prohibición de \gls{doble-reserva}, que se expresa como \gls{indice-filtrado} sobre los envíos abiertos, y la protección del último administrador, que necesita serializar un tramo crítico mediante un cerrojo de aviso. MySQL no ofrece índices parciales y habría obligado a simular la condición con una columna derivada; SQLite no admite escritores concurrentes, que es precisamente aquello a lo que esas dos reglas deben resistir; y SQL Server sí dispone de índices filtrados, pero añade una decisión de licencia y una imagen mayor sin aportar nada que el proyecto necesite. PostgreSQL reúne ambas capacidades, cuenta con proveedor EF Core mantenido e imagen oficial, y no impone condiciones de uso.

La autenticación utiliza \gls{jwt}, definido por RFC~7519 \cite{rfc7519}. El token firmado transporta identidad y rol, y evita mantener un almacén de sesiones que habría añadido un componente al despliegue y un estado que conservar entre reinicios. La alternativa ---sesión en servidor con un identificador opaco en cookie--- ofrece revocación inmediata por diseño, y conviene decir que el endurecimiento acabó recuperando esa propiedad sin adoptar el mecanismo: el token viaja en una cookie \texttt{HttpOnly} y una versión de sesión comprobada en cada petición lo invalida al instante. La elección final es por tanto un punto intermedio deliberado y no una API sin estado en sentido estricto. ASP.NET Core Identity aporta el \emph{hasher} adaptativo de contraseñas; no se usa el modelo completo de Identity porque dos roles, \gls{bootstrap} y administración controlada no justifican sus tablas y flujos adicionales.

\section{Frontend}

El frontend es una \gls{spa} con React 19 \cite{react}. El modelo de componentes permite construir pantallas de catálogo, formularios y detalle operativo compartiendo piezas de presentación sin mezclar reglas del servidor. Se escogió una SPA frente a vistas renderizadas por la API porque el flujo de asignación, filtros en URL y actualización de cronología se benefician de interacción local; la separación también hace explícito el contrato de integración que el TFG pretende documentar.

Entre las bibliotecas que resuelven ese modelo, la decisión atiende a dos razones concretas. Angular habría traído inyección de dependencias, módulos y programación reactiva propios, una superficie de framework mayor de la que necesitan las pantallas de este producto, y con ella una parte del tiempo dedicada a aprender el marco en lugar del problema. Vue y Svelte lo resolverían igual de bien, pero React tiene detrás la herramienta de prueba con la que este trabajo quería verificar el frontend: consultar la interfaz por rol accesible y texto visible en lugar de por estructura interna. Dicho de otro modo, la biblioteca se eligió por su ecosistema de verificación y no por sus propiedades de renderizado, que para cinco pantallas son equivalentes.

TypeScript añade comprobación estática sobre propiedades, respuestas y estados de interfaz \cite{typescript}. Su coste de anotación se compensa especialmente en un cliente que consume numerosos \glspl{dto}; un cambio de contrato puede detectarse durante la compilación antes de convertirse en una pantalla rota. El frontend no replica las entidades C\# automáticamente: mantiene tipos orientados a su uso y conserva la API como frontera.

Vite proporciona servidor de desarrollo, proxy y compilación de producción \cite{vite}. Frente a configuraciones manuales de \emph{bundling}, ofrece un punto de partida pequeño y tiempos de realimentación breves. Frente a un armazón con opinión propia sobre el servidor, como los que ofrecen renderizado en servidor y rutas de backend, se descartó porque este producto ya tiene su servidor y su contrato: añadir un segundo lugar donde pueda vivir lógica habría difuminado justamente la frontera que la arquitectura pretende hacer visible.

React Router define las URL, los parámetros, las rutas protegidas y la navegación \cite{reactrouter}. Los filtros de envíos se representan en la \emph{query string}, lo que aprovecha el navegador como parte del estado de navegación. Se prefirió a resolver el enrutado sobre la API de historial del navegador porque las tres cosas que el producto necesita ---rutas cuya autorización depende del rol, un contenedor común de navegación y filtros que viven en la dirección--- son exactamente las que la biblioteca ya resuelve, y escribirlas a mano habría añadido código propio que mantener sin comportamiento nuevo que mostrar.

El cliente usa \gls{json} sobre HTTP y centraliza token, sobres de éxito y errores. No se incorporó una biblioteca global de estado: contexto para autenticación y estado local por pantalla cubren el alcance actual con menos acoplamiento. Tampoco se añadió un sistema visual externo; CSS propio permite mantener una interfaz coherente y accesible sin depender de componentes cuya personalización excedería el beneficio para este producto.

\section{Pruebas y calidad}

xUnit verifica el backend \cite{xunit}. Frente a NUnit o MSTest las diferencias funcionales son menores, y pesó una propiedad concreta: xUnit crea una instancia de la clase por cada prueba, lo que encaja con el patrón que esta suite necesita ---un contexto de base de datos nuevo y aislado en cada caso--- sin recurrir a código de preparación compartido capaz de filtrar estado de una prueba a la siguiente. Las pruebas de servicio ejercitan reglas y persistencia con una infraestructura aislada; las de controlador comprueban políticas, códigos y contrato. Esta separación localiza fallos mejor que probarlo todo exclusivamente a través de HTTP, pero mantiene cobertura del borde real. Las migraciones se validan además contra PostgreSQL en \gls{ci}, porque un proveedor en memoria no reproduce índices, tipos ni restricciones del motor elegido.

Vitest comparte el modelo de módulos y configuración de Vite \cite{vitest}; React Testing Library consulta la interfaz por roles, etiquetas y texto visible. Se evita acoplar las pruebas a detalles internos de componentes: una prueba debe observar lo que una persona puede hacer, como filtrar, recibir un aviso o perder acceso a una ruta. Oxlint y la compilación de TypeScript complementan las suites con análisis estático. Se prefirió Oxlint a ESLint, la opción habitual en este ecosistema, porque cubre el conjunto de reglas que este proyecto necesita sin requerir configuración propia y con un tiempo de ejecución que permite invocarlo en cada cierre de sprint sin convertirlo en un trámite que se acabe omitiendo.

\section{Ejecución y automatización}

Docker crea imágenes separadas para API y frontend, y reutiliza la imagen oficial de PostgreSQL \cite{docker}. Docker Compose describe tres servicios, dependencias, comprobación de salud, variables y volumen. La alternativa consistía en pedir a quien reproduzca el trabajo que instale .NET, Node y PostgreSQL en su equipo y los configure siguiendo una lista de pasos; se descartó porque haría depender de esa máquina una afirmación central de la memoria, la de que el sistema se levanta desde un documento. Un fichero de composición convierte esa afirmación en una orden única y comprobable, y es también lo que permite que el entorno accesible del Capítulo~\ref{chap:despliegue} reutilice la misma descripción con otros parámetros en lugar de mantener una paralela que divergiría.

El contenedor web usa Nginx tanto para servir los estáticos compilados como para reenviar \texttt{/api}. Habría sido posible ahorrar ese contenedor sirviendo la SPA desde la propia API, pero entonces el desarrollo ---donde el proxy es Vite--- y la ejecución contenerizada tendrían fronteras distintas, y el frontend dejaría de poder construirse y publicarse sin reconstruir también el backend. Con Nginx, una demostración local atraviesa la misma frontera que el desarrollo y que el despliegue. La configuración sensible se obtiene de un \texttt{.env} ignorado y el Compose falla de forma explícita si falta una variable obligatoria.

GitHub Actions ejecuta restauración, \emph{lint}, compilación, pruebas y validación de migraciones \cite{githubactions}. Se eligió por proximidad y no por preferencia: el repositorio ya está alojado en GitHub, de modo que la integración no añade ninguna cuenta, ningún servicio externo ni ningún secreto que custodiar, y su cuota gratuita para repositorios públicos cubre el uso del proyecto. Un servidor propio, del tipo Jenkins, habría exigido operarlo, que es precisamente la actividad que este trabajo declara fuera de su alcance; y una plataforma de terceros habría introducido una credencial de acceso al código a cambio de capacidades que aquí no se emplean. El mismo argumento explica que el registro de contenedores del que se sirve el despliegue sea el de la propia cuenta.

La automatización no constituye un despliegue complejo: es una red de seguridad que demuestra que el repositorio puede reconstruirse fuera del equipo de desarrollo. Esta realimentación frecuente se alinea con la entrega continua \cite{humble2010}. El destino del despliegue se decidió deliberadamente en el séptimo incremento y no antes: escogerlo sin conocer las necesidades reales del producto habría supuesto anteponer la infraestructura a lo que debe sostener.

Vista en conjunto, la selección comparte una propiedad que la estimación de costes recoge más adelante: ninguna de estas herramientas exige licencia ni cuota para este uso. No fue una restricción impuesta desde fuera, sino un criterio de la propia elección, porque un trabajo que solo pudiera reproducirse pagando dejaría de ser reproducible para quien lo lee.
